\begin{frame}[fragile]{리뷰 템플릿 --- 빈칸이 채워지지 않으면 3점 미만}
\begin{lstlisting}
팀 이름: ________   리뷰어: ________
1. 문제 정의 -- 무엇을 예측/분류하려 했는지 한 문장. 명확했는가?
2. 방법론 적절성 -- 데이터의 어떤 특성(불균형, 차원, 결측치)에
   어떤 모델이 맞았는가? 그 선택에 동의하는가, 근거와 함께.
3. 수식적 정당화 -- "최소 1개"(8.1)를 채웠는가? 어떤 장의 개념을
   썼고, 그 설명이 *실제로* 타당한가?
4. 결과의 정직성 -- 잘 안 된 부분과 그 원인을 분석했는가?
5. 반박 가능한 질문 (필수, 1개 이상)
   근거 개념: Ch__ 절__ (개념: ______)
   질문: ________ (답이 나쁘면 결과가 흔들려야 한다)
\end{lstlisting}
  \vskip0.2em
  {\footnotesize 5번의 빈칸이 채워지지 않으면 리뷰는 1$\sim$2점에
  머문다 --- ``개념''과 ``질문'' 두 칸이 모두 비어 있어야
  \textbf{반박 불가능한} 리뷰다. 템플릿 채우기가 곧 ``구체적인가?''
  라는 채점 기준의 자기 점검.}
\end{frame}
